\unit{강연 232 — 피카르 스킴: 존재 정리}
\sid{232}
\src{143}
\seminar{제14년차, 1961/62, 제232호\\1962년 2월}
\paperhead{대수기하학에서의 하강 기법과 존재 정리\\
V. 피카르 스킴: 존재 정리}%
  {알렉상드르 그로텐디크 지음}

\fsection{1}{상대 피카르 군과 함자}

모든 준스킴(더 일반적으로는 모든 환 달린 공간) $X$에 대하여, $X$ 위의
가역 가군(즉 $\mathcal O_X$와 국소적으로 동형인 가군)들의 동형류로
이루어진 군을 $X$의 (절대) \emph{피카르 군}이라 하고
$\operatorname{Pic}(X)$로 나타낸다. 따라서 표준 동형사상
\begin{equation}
\operatorname{Pic}(X)\xrightarrow{\sim}H^1(X,\mathcal O_X^*)
\tag{1.1}
\end{equation}
이 있다. 여기서 $\mathcal O_X^*$는 $\mathcal O_X$의 가역원층을 나타낸다.
(실제로 이는 표준 가역 가군 $\mathcal O_X$의 자기동형사상들로 이루어진
층과 동일시된다.) 분명히 $X\longmapsto\operatorname{Pic}(X)$는 $X$에 관한
반변 함자이며, 동형사상 (1.1)은 함자적이다.

$X$가 다른 준스킴 $S$ 위의 준스킴이라고 하자. $S$ 위의 준스킴들의 범주
$(\mathrm{Sch})/S$에서 $S'$이 움직일 때, 앞의 논의에 따라
\[
S'\longmapsto\operatorname{Pic}(X\times_S S')
\]
라는 반변 함자를 얻는다. 분류하려는 가역 가군들이 자명하지 않은
자기동형사상을 가지므로 이 함자는 “국소적 성격”을 갖지 않으며([5], IV,
5.4), 따라서 “표현 가능”할 가능성이 전혀 없다([4], II, A). 그러므로 이
함자를 “국소적으로 만들어야” 한다. 일반적으로 모든 상대 준스킴 $X/S$에
대하여 상대적인 군
\begin{equation}
\operatorname{Pic}'(X/S)=H^0\bigl(S,R^1f_*(\mathcal O_X^*)\bigr)
\tag{1.2}
\end{equation}
을 도입하자. 여기서 $f:X\longrightarrow S$는 구조사상이다([4], II, C 3과
비교하라). 인용한 곳에서는 이 군을 상대 피카르 군이라 부르지만, 곧
드러날 이유로 여기서는 $X/S$의 \emph{제한 상대 피카르 군}이라 부르는
편이 낫다. $S'$이 $(\mathrm{Sch})/S$에서 움직이면
\[
S'\longmapsto\operatorname{Pic}'(X\times_S S'/S')
\]
은 $S'$에 관한 반변 함자이다. 이 함자도
$\operatorname{Pic}'_{X/S}$로 나타내며, 본질적으로 다음 식으로 정의된다.
\begin{equation}
\operatorname{Pic}'_{X/S}(S')=
\operatorname{Pic}'(X\times_S S'/S').
\tag{1.3}
\end{equation}

\src{144}
필요한 조치를 취했으므로 이제 이 함자는 “국소적 성격”을 갖는다.
직관적으로 (1.3)의 우변은 매개변수 준스킴 $S'/S$로 지표화된, $X/S$의
(섬유들 위의) 가역층 동형류들의 “대수적 족”의 집합으로 해석된다. 함자
$\operatorname{Pic}'$가 표현 가능할 때, 이를 표현하는 $S$ 위의 준스킴을
$\operatorname{Pic}'_{X/S}$로 나타내고 $S$ 위의 $X$의
\emph{피카르 준스킴}이라 한다. 따라서 이 경우
\begin{equation}
\operatorname{Hom}_S(S',\operatorname{Pic}'_{X/S})
\simeq\operatorname{Pic}'_{X/S}(S')
=\operatorname{Pic}'(X\times_S S'/S')
\tag{1.4}
\end{equation}
을 갖는다.

그러나 $\operatorname{Pic}'_{X/S}$가 표현 가능하지 않은데도 상대 피카르
준스킴을 자연스럽게 정의할 수 있는 중요한 경우들이 있다. 예를 들어 체
$k$ 위에 $k$-유리점이 없는 “브라우어--세베리” 다양체가 그러하다. 이는
절대 피카르 군 $\operatorname{Pic}(X\times_S S'/S')$으로부터
$\operatorname{Pic}'$를 정의할 때 아직 충분히 국소화하지 않았기
때문이다. 정확히 말하면 $\operatorname{Pic}'$는 일반적으로 “충실 평탄
하강과 양립”하지 않는다. 이를 자세히 설명하자.

$\mathfrak M$을 \emph{충실 평탄이고 준콤팩트인} 준스킴 사상들의 모임이라
하자. 이 모임은 밑변경과 합성에 대해 안정하다. $P$를
$(\mathrm{Sch})/S$에서 집합의 범주로 가는 반변 함자라 하고, 모든
$S$-사상 $u:T'\longrightarrow T$, $u\in\mathfrak M$에 대하여 도식
\begin{equation}
P(T)\longrightarrow P(T')\rightrightarrows P(T'\times_TT')
\tag{1.5}
\end{equation}
을 생각하자. 이는 도식
\[
T\longleftarrow T'\mathrel{\substack{\xleftarrow{\,\mathrm{pr}_1\,}\\[-0.35ex]
\xleftarrow[\,\mathrm{pr}_2\,]{}}}T'\times_TT'
\]
에 $P$를 적용하여 얻는다. $P$가 표현 가능하면 하강 이론([4], I, B,
정리 2)에 따라 모든 $u\in\mathfrak M$에 대하여 도식 (1.5)가 완전하다.
이를 $P$가 $\mathfrak M$과 양립한다고 말한다. 지금의 경우에는 $P$가
“충실 평탄 하강과 양립”한다고도 하고, $(\mathrm{Sch})/S$ 위의 “준층”
$P$가 $\mathfrak M$이 주는 국소화 개념에 관한 “층”이라고도 한다. $P$가
임의의 함자일 때, 보통의 위상적 국소화에서 잘 알려진 표준 절차를 통해
“층” $\mathcal P$와 함자 준동형 $P\longrightarrow\mathcal P$를 대응시킬
수 있으며, 이 준동형은 자명한 의미에서 보편적이다. $\mathcal P$는 다음과
같이 구체적으로 계산할 수 있다. $\mathcal P(T)$를 정의하기 위하여,
$u:T'\longrightarrow T$가 $\mathfrak M$에 속하는 모든 $T$ 위의 $T'$에
대해 다음 집합을 생각한다.

\src{145}
$\overline H^0(T'/T,P)$를 $P(T')$의 원소 $\xi$ 가운데, 그 두 상
$\xi_1,\xi_2\in P(T'\times_TT')$가 다음 성질을 갖는 것들로 이루어진
부분집합이라 하자. 즉 $v\in\mathfrak M$인 사상
$v:T''\longrightarrow T'\times_TT'$가 존재하여, $\xi_1$과 $\xi_2$의
$P(T'')$에서의 상이 같아야 한다. [주의. 이와 같이 정의한 집합
$\overline H^0$는 [4], I, A, 4(a)에서 도입한 집합 $H^0(T'/T,P)$보다
크다.] $T$를 고정하고 $T'$을 움직이면(여전히 $u\in\mathfrak M$이라
하고, $T'$들의 모임에는 지배관계가 정하는 준순서를 준다)
$\overline H^0(T'/T,P)$들은 귀납계를 이루며, 다음과 같이 놓는다.
\begin{equation}
\mathcal P(T)=\varinjlim_{T'}\overline H^0(T'/T,P).
\tag{1.6}
\end{equation}
이 식의 $T$에 관한 함자 법칙은 분명하게 구체화된다.

\[
P(T)=\operatorname{Pic}(X\times_ST)
\]
일 때, (1.6)이 정의하는 $(\mathrm{Sch})/S$ 위의 반변 함자를 $S$ 위의
$X$의 \emph{상대 피카르 함자}라 하고 $\operatorname{Pic}_{X/S}$로
나타낸다. 또한 군 $\operatorname{Pic}_{X/S}(S)$를 $S$ 위의 $X$의
\emph{상대 피카르 군}이라 하고 $\operatorname{Pic}(X/S)$로 나타낸다.
그러면 분명한 전단사
\begin{equation}
\operatorname{Pic}_{X/S}(T)\simeq\operatorname{Pic}(X\times_ST/T)
\tag{1.7}
\end{equation}
를 얻는다.

따라서 $\operatorname{Pic}(X/S)$의 한 원소는 다음과 같은 자료로
정의된다. 충실 평탄이고 준콤팩트인 사상 $S'\longrightarrow S$와 군
$\operatorname{Pic}(X\times_SS')$의 원소 $\xi'$를 택하되, 충실 평탄이고
준콤팩트인 사상 $S''\longrightarrow S'\times_SS'$가 존재하여 $\xi'$의
두 역상이 $\operatorname{Pic}(X\times_SS'')$에서 같아야 한다. 방금
설명한 조건을 만족하는
\[
\xi'\in\operatorname{Pic}(X\times_SS'),\qquad
\xi_1\in\operatorname{Pic}(X\times_SS_1)
\]
이 $\operatorname{Pic}(X/S)$의 같은 원소를 정의할 필요충분조건은,
충실 평탄이고 준콤팩트인 사상
$S'_1\longrightarrow S'\times_SS_1$이 존재하여 문제의 두 원소의 상이
$\operatorname{Pic}(X\times_SS'_1)$에서 같은 것이다. 앞에서 도입한 함자
$P'=\operatorname{Pic}'_{X/S}$로 계속 작업하는 편이 흔히 편리하다.
표준 사상 $P\longrightarrow P'$이 동형사상
\begin{equation}
\mathcal P\xrightarrow{\sim}\mathcal P'
\tag{1.8}
\end{equation}
을 유도함을 곧바로 알 수 있다. 이는
$\operatorname{Pic}'_{X/S}=P'$를 이용하여
$\operatorname{Pic}_{X/S}$를 기술하는, 흔히 더 편리한 방법을 준다. 아래
2.3에 따라, 방금 설명한 $\operatorname{Pic}(X/S)$의 기술에서 $P$를
$P'$로 바꾸면 실제로

\src{146}
$S''=S'\times_SS'$ 및 $S'_1=S'\times_SS_1$로 잡을 수 있다. 적어도 인용한
곳에 명시된 조건 아래에서는 그렇다.

함자 $\operatorname{Pic}_{X/S}$가 표현 가능하면 $X/S$가 피카르 준스킴을
갖는다고 말한다. 이 함자를 표현하는 $S$ 위의 준스킴을 $S$ 위의 $X$의
\emph{피카르 준스킴}이라 하고, 여전히 $\operatorname{Pic}_{X/S}$로
나타낸다. 이를 위해서는 $P'=\operatorname{Pic}'_{X/S}$가 표현
가능하기만 하면 됨이 분명하다. 실제로 이 경우 $P'$는 이미 “층”이고,
(1.8)에 따라 $P'\longrightarrow\mathcal P'$은 표준 사상
\begin{equation}
\operatorname{Pic}'_{X/S}\longrightarrow\operatorname{Pic}_{X/S}
\tag{1.9}
\end{equation}
과 동일시되며, 이 사상은 동형사상이다. 따라서 우리의 용어는 (1.4)와
함께 앞에서 도입한 용어와 양립한다. 일반적으로
$\operatorname{Pic}_{X/S}$가 존재할 때, 이는 함자적 동형사상
\begin{equation}
\operatorname{Hom}_S(S',\operatorname{Pic}_{X/S})\xrightarrow{\sim}
\operatorname{Pic}(X\times_SS'/S')
\tag{1.10}
\end{equation}
으로 정의된다.

\fsection{2}{여러 상대 및 절대 피카르 군 사이의 관계}

\result{명제}{2.1} $f:X\longrightarrow S$가
$\mathcal O_S\xrightarrow{\sim}f_*(\mathcal O_X)$를 만족하는 사상이라고
하자. 그러면 완전열
\[
0\longrightarrow\operatorname{Pic}(S)\longrightarrow\operatorname{Pic}(X)
\longrightarrow\operatorname{Pic}'(X/S)
\]
이 있다. $X$가 $S$ 위의 단면을 가지면 마지막 사상은 전사이다. 즉
동형사상
\[
\operatorname{Pic}'(X/S)\xrightarrow{\sim}
\operatorname{Pic}(X)/\operatorname{Pic}(S)
\]
이 있다.

이 완전열은 $f$와 $\mathcal O_X^*$에 관한 르레 스펙트럴 열에 대응하는
저차 완전열로 볼 수 있다. 두 번째 주장도 마찬가지로 형식적이다.

\result{명제}{2.2} $f:X\longrightarrow S$가 준콤팩트 분리 사상이고
$\mathcal O_S\xrightarrow{\sim}f_*(\mathcal O_X)$를 만족한다고 하자. 또한
$S'\longrightarrow S$를 충실 평탄이고 준콤팩트인 사상이라 하자. 그러면
다음이 성립한다.

\noindent (i) $\operatorname{Pic}'(X/S)\longrightarrow
\operatorname{Pic}'(X\times_SS'/S')$은 단사이다.

\noindent (ii) $X$가 $S$ 위에서 국소적으로 단면을 갖는다고 하자. 즉 모든
$s\in S$가 열린 근방 $U$를 가져서 $X|U$가 $U$ 위의 단면을 갖는다고
하자. 그러면 도식
\[
\operatorname{Pic}'(X/S)\longrightarrow
\operatorname{Pic}'(X\times_SS'/S')\rightrightarrows
\operatorname{Pic}'(X\times_SS''/S'')
\]

\src{147}
은 완전하다. 여기서 $S''=S'\times_SS'$이다.

첫 번째 주장은 충실 평탄 하강의 초등적 성질과 다음의 일반적 관찰로부터
따라온다. $f:X\longrightarrow S$가
$\mathcal O_S\xrightarrow{\sim}f_*(\mathcal O_X)$를 만족하는 사상이라
하자. $S$ 위의 유한 계수 국소 자유 가군들의 범주에서 $X$ 위의 유한 계수
국소 자유 가군들의 범주로 가는 함자
\[
\mathcal F\longmapsto f^*(\mathcal F)
\]
는 충만하고 충실하다. 그 본질적 상은 $X$ 위의 가군 $\mathcal G$ 가운데
$f_*(\mathcal G)$가 국소 자유이고 표준 준동형
\[
f^*\bigl(f_*(\mathcal G)\bigr)\longrightarrow\mathcal G
\]
이 동형사상인 것들로 이루어진다. 두 번째 주장은 [4], I, B, 4에서 하강
이론으로 증명되었다.

2.2의 결과는 다음과 같이도 진술할 수 있다.

\result{따름정리}{2.3} 2.2의 조건 아래에서 표준 준동형 (1.9)
\[
\operatorname{Pic}'_{X/S}\longrightarrow\operatorname{Pic}_{X/S}
\]
은 단사이고, $X$가 $S$ 위에서 국소적으로 단면을 가지면 전단사이기까지
하다. (따라서 후자의 경우 상대 피카르 군 $\operatorname{Pic}(X/S)$과
제한 상대 피카르 군 $\operatorname{Pic}'(X/S)$은 동일하다.)

이를 2.1과 결합하면 다음을 얻는다.

\result{따름정리}{2.4} 2.2의 조건 아래에서 완전열
\[
0\longrightarrow\operatorname{Pic}(S)\longrightarrow\operatorname{Pic}(X)
\longrightarrow\operatorname{Pic}(X/S)
\]
이 있다. $X$가 $S$ 위의 단면을 가지면 마지막 준동형은 전사이다. 즉
동형사상
\[
\operatorname{Pic}(X/S)\xrightarrow{\sim}
\operatorname{Pic}(X)/\operatorname{Pic}(S)
\]
이 있다.

\result{비고}{2.5} $f:X\longrightarrow S$가
$\mathcal O_S\xrightarrow{\sim}f_*(\mathcal O_X)$를 만족하는 사상이고,
$g$가 $S$ 위의 $X$의 단면이라고 하자. $\mathcal L$을 $X$ 위의 가역
가군이라 하자. 동형사상
\[
\mathcal O_S\xrightarrow{\sim}g^*(\mathcal L)
\]
을 $\mathcal L$의 $g$-강체화라 하고, $g$-강체화가 주어진 $X$ 위의 가역
가군 $\mathcal L$을 \emph{$g$-강체화된 가역 가군}이라 하자. 이러한
구조의 모든 자기동형사상은 자명하며, $\operatorname{Pic}'(X/S)$은 $X$
위의 $g$-강체화된 가역 가군들의 동형류 집합과
동일시된다.\footnote{인쇄본에는 “$S$ 위”라고 되어 있다. 그러나 단면
$g:S\to X$에 의한 강체화는 여기서 $X$ 위의 가군 $\mathcal L$에 관한
것이다.} (바로 이 사실 덕분에 2.2(ii)를 증명할 때 하강 이론을 사용할 수
있었다.) 이는 적어도 $f$가 더 나아가 준콤팩트이고 분리인 경우
$\operatorname{Pic}(X/S)$에 대한 새로운 해석을 준다. 실제로 이 경우
2.3에 따라
\[
\operatorname{Pic}(X/S)\simeq\operatorname{Pic}'(X/S)
\]
이다.

\src{148}
\result{비고}{2.6} $f:X\longrightarrow S$를 2.2에서와 같은 사상이라
하고, $S'\longrightarrow S$를 다음 조건을 만족하는 충실 평탄이고
준콤팩트인 사상이라 하자. 즉 $S$-사상 $S'\longrightarrow X$가 존재하며,
동치로 $X'=X\times_SS'$가 $S'$ 위의 단면을 갖는다고 하자.
$S''=S'\times_SS'$, $X''=X\times_SS''$로 놓고 완전열
\[
\operatorname{Pic}(X/S)\longrightarrow\operatorname{Pic}(X'/S')
\rightrightarrows\operatorname{Pic}(X''/S'')
\]
을 생각하자. 2.4를 적용하면 완전열
\[
\operatorname{Pic}(X/S)\longrightarrow
\operatorname{Pic}(X')/\operatorname{Pic}(S')
\rightrightarrows
\operatorname{Pic}(X'')/\operatorname{Pic}(S'')
\]
을 얻는다. 특히 상대 피카르 군의 모든 원소는 이미
$\operatorname{Pic}(X')$의 어떤 원소에서 “온다”. 따라서 이는 앞
번호에서 설명한 상대 피카르 군의 기술을 크게 단순화하며, $S$ 위의 $X$의
피카르 함자에 대해서도 마찬가지다. 실제로 $S$ 위의 모든 $T$에 대하여
$X\times_ST/T$와 사상 $T'=S'\times_ST\longrightarrow T$에 앞의 논의를
적용할 수 있다. 예를 들어 $f$ 자체가 충실 평탄이면 $S'=X$로 잡을 수
있다. 더 나아가 $f$가 유한형(각각 매끄러운 등)이면, 상대 피카르 함자
\[
T\longmapsto\operatorname{Pic}(X\times_ST/T)
\]
를 기술할 때 유한형(각각 매끄러운 등)인 밑변경
$T'\longrightarrow T$만을 사용해도 된다. $f$가 사영이고 평탄하며 $S$가
국소 뇌터일 때에는, 위의 $S'\longrightarrow S$를 다음과 같이 잡을 수
있음이 증명된다. $S$를 덮는 열린집합 $S_i$들의 평탄 피복 $S'_i$를 택해
$S'$이 이들의 서로소 합이 되게 할 수 있다. $f$가 더 나아가 분리
가능하면 $S'_i$를 $S_i$ 위에서 에탈이 되게 잡을 수 있다.

\fsection{3}{주요 존재 정리: 진술}

알려진 모든 경우를 포괄하는 피카르 준스킴의 존재 명제는 추측의 형태로도
아직 없다. 말하자면 “실질적으로 필요한” 조건은
$f:X\longrightarrow S$가 고유이고(이는 본질적인 유한성 성질들을
보장한다) 평탄인 것이다. 그러나 $S$가 체 $k$ 위의 쌍대수
$k[t]/(t^2)$의 스펙트럼이고(가령 $k$를 복소수체 $\mathbf C$로 잡아도
된다) $X$가 1차원인 경우에도 이 조건들은 충분하지 않다. 이 강연을 쓰는
현재, 피카르 준스킴에 관한 가장 중요한 존재 정리들은 다음 정리에서
따라온다.

\result{정리}{3.1} $f:X\longrightarrow S$를 국소 뇌터 준스킴들의
사상이라 하자. 다음을 가정한다.

\noindent (i) $f$는 사영이다.

\src{149}
\noindent (ii) $f$는 평탄이다.

\noindent (iii) $f$의 기하적 섬유들은 정역 준스킴이다.

이 조건들 아래에서 $\operatorname{Pic}_{X/S}$가 존재한다.

뒤의 두 번호에서 개략적으로 제시할 증명은 동시에 다음도 보일 것이다.
$\xi$를 $S$에 관해 매우 풍부한 층 $\mathcal O_X(1)$에 대응하는
$\operatorname{Pic}_{X/S}$의 단면이라 하자. (즉 이 층은 사영 몰입 사상
$X\longrightarrow\mathbf P(\mathcal E)$에 의해 유도된다.) 그러면
$\operatorname{Pic}_{X/S}$의 열린 부분 $U$가 존재하여, $U$는 $S$ 위에서
준사영인 열린 부분들의 서로소 합이고 $\xi$에 의한 평행이동에 대해
안정하며, $\operatorname{Pic}_{X/S}$는 열린집합 $U-n\xi$들의 증가하는
합집합이다. (이 열린집합들은 모두 $U$와 동형이다.) 특히 3.1의 조건
아래에서 $\operatorname{Pic}_{X/S}$는 $S$ 위에서 분리되어 있다.

\result{비고}{3.2} 예들을 통해 다음을 알 수 있다. (예를 들어 $S$가 이산
값매김환의 스펙트럼이고 $X$가 $S$ 위에서 상대 차원 1인 경우를 생각할 수
있다.) 3.1에서 가정 (iii)을 생략하고, 모든 $s\in S$에 대하여 준동형
\[
k(s)\longrightarrow H^0(X_s,\mathcal O_{X_s})
\]
이 동형이라는 더 약한 가정으로 바꾸면 $\operatorname{Pic}_{X/S}$는
$S$ 위에서 분리되어 있지 않을 수도 있다. 이는 $f$의 기하적 섬유들이
축소이지만 정역인 기하적 일반섬유 하나가 특수화될 때 두 기약 성분으로
``갈라지는'' 경우에도, 기하적 섬유들이 기약이지만 정역인 기하적
일반섬유 하나가 ``다중 섬유''로 특수화되는 경우에도 일어난다. 첫 번째
경우는 예를 들어 원뿔곡선이 한 점에서 만나는 두 직선으로 퇴화할 때
나타난다. 두 번째 경우의 예는 D. 멈퍼드가 알려 주었는데, 타원곡선이
이중 타원곡선으로 퇴화하는 예이다. 이 예들은 임의의 표수에서 성립한다.

\result{비고}{3.3} 3.1의 조건 아래에서 $\operatorname{Pic}_{X/S}$가
$S$ 위에서 유한형이고 따라서 준사영인 열린 부분들의 서로소 합인지는
나는 알지 못한다.\footnote{1962년 정오표는 D. 멈퍼드가 여기서 전개한
피카르 스킴 이론을 이용하여 이 문제를 긍정적으로 해결했다고 전한다.}
힐베르트 스킴의 경우와 마찬가지로([4], IV), 힐베르트 다항식
$Q\in\mathbf Q[t]$를 생각하면 $\operatorname{Pic}_{X/S}$를 열린집합
$\operatorname{Pic}^{Q}_{X/S}$들의 서로소 합으로 분해할 수 있다. 이
열린집합들은 $S$ 위에서 준사영일 것으로 보이며, 적어도 다음 강연에서
$f$가 매끄러운 사상일 때 이를 보게 될 것이다. 가정 (i)을 ``$X$는
$S$ 위에서 국소적으로 사영이다''라는 가정으로 바꾸는 경우에는 주의해야
한다. ($\operatorname{Pic}_{X/S}$의 존재 문제는 분명히 $S$ 위에서
국소적이므로, 이 가정만으로도 3.1은 성립한다.) 반면 이 경우에는

\src{150}
$\operatorname{Pic}_{X/S}$가 $S$ 위에서 유한형이 아닌 연결 성분들을
포함하는 예를 쉽게 들 수 있다. 예를 들어 $X_0$를 대수적으로 닫힌 체
$k$ 위의 비특이 사영 대수다양체라 하고, $u$를 그 자동사상, $\xi$를
$X_0$의 네론--세베리 군의 원소라 하되 $u^n(\xi)$들이 서로 다르다고
하자. 이를테면 $X_0$로 타원곡선 $E$와 자기 자신의 곱을 취하고, $u$로는
$E\times E$의 자동사상
\[
(x,y)\longmapsto(x,y+x)
\]
을 취할 수 있다. $S$를 두 점 $a$와 $b$에서 만나는 두 비특이 기약곡선의
합집합이라 하자. $S$ 위에는 연결 $\mathbf Z$-토서 $P$가 있다.
$u$가 정하는 $\mathbf Z$의 $X_0$ 위의 작용들을 이용하면, 섬유가 $X_0$인
$S$ 위의 연관 다발을 얻는다. 이 다발은 $S-a$와 $S-b$ 위에서 자명하고,
지금 생각하는 특별한 경우에는 실제로 $S$ 위의 아벨 스킴이다. 또한
$\operatorname{Pic}_{X/S}$는 $P$와, $u$를 통해
$\operatorname{Pic}_{X_0/k}$ 위에 주어지는 $\mathbf Z$의 작용들에
연관된 다발이며, 다음과 동형인 연결 성분을 포함함을 쉽게 알 수 있다.
\[
P\times\operatorname{Pic}^{0}_{X_0/k}.
\]
(여기서 $\operatorname{Pic}^{0}$는 $\operatorname{Pic}$ 안에서 항등원의
연결 성분을 나타낸다.) 이 성분은 $S$ 위에서 유한형이 아니다. [$S$ 위에서
분리되지 않은 피카르 준스킴의 여러 경우에도 3.2에서 생각한 것과 비슷한
현상이 일어난다.]

\fsection{4}{상대 카르티에 인자와 사영 다발}

우리는 양의 인자만 사용할 것이며, 이 번호의 나머지 부분에서는 이
덧붙인 수식어를 생략한다.

$X$를 준스킴이라 하자. $X$ 위의 \emph{카르티에 인자}, 또는 간단히
$X$ 위의 인자란 가역 가군인 아이디얼 $\mathcal J$에 의해 정의되는
$X$의 닫힌 부분준스킴 $D$를 말한다. 즉 $\mathcal J$는
$\mathcal O_X$의 영인자가 아닌 한 절단에 의해 국소적으로 생성된다.
$D$에 가역 가군
\[
\mathcal L(D)=\mathcal J^{-1}
\]
을 결부한다. 그리고 표준 포함 사상
$\mathcal J\longrightarrow\mathcal O_X$는 표준 준동형
\[
s_D:\mathcal O_X\longrightarrow\mathcal J^{-1}=\mathcal L(D),
\qquad\text{즉}\quad s_D\in\Gamma(X,\mathcal L(D))
\]
을 준다. 더 나아가 인자 하나를 주는 것은 본질적으로 $X$ 위의 가역 가군
$\mathcal L$과 모든 점에서 영인자가 아닌 절단 $s$를 주는 것과 동치이다.
역으로 이러한 쌍 $(\mathcal L,s)$에는 $s$의 ``인자''를 결부하며, 이를
$\operatorname{div}(s)$로 나타낸다. $X$ 위의 주어진 가역 가군
$\mathcal L$을 정의하는 인자 $D$들의 집합은 몫집합

\src{151}
\[
\Gamma(X,\mathcal L)^*/\Gamma(X,\mathcal O_X^*)
\]
과 전단사로 대응한다. 여기서 $\Gamma(X,\mathcal L)^*$는
$\Gamma(X,\mathcal L)$의 절단들 가운데 모든 점에서 영인자가 아닌 것들로
이루어진 부분집합을 나타낸다.

이제 국소 유한형 사상 $f:X\longrightarrow S$가 주어졌다고 하고, 간단히
$S$가 국소 뇌터라고 가정하자. $\mathcal J$를 $X$ 위의 연접 아이디얼층,
$D$를 이것이 정의하는 $X$의 부분스킴이라 하고, $x\in X$ 및 $s=f(x)$라
하자. 다음 조건들은 서로 동치이다.

\noindent (i) $\mathcal J$는 $x$에서 가역이고(즉 $\mathcal J_x$는
$\mathcal O_{X,x}$의 정칙 원소 하나로 생성된다), $D$는 $x$에서
$S$ 위로 평탄이다.

\noindent (ii) $X$와 $D$는 $x$에서 $S$ 위로 평탄이고, $D_s$는 점
$x$에서 섬유 $X_s$ 위의 카르티에 인자이다.

\noindent (iii) $X$는 $x$에서 $S$ 위로 평탄이고, $\mathcal J_x$는
$X_s$ 위에 영인자가 아닌 싹을 유도하는 원소 $f_x$로 생성된다.

이때 $D$를 생각하는 점에서 $X/S$ 위의 \emph{상대 카르티에 인자}(또는
간단히 $X/S$ 위의 상대 인자)라고 한다. (i)에서 $D$가 $x$의 근방에 있는
점들에서도 상대 인자임을 알 수 있다. 따라서 $X$와 $D$가 $S$ 위에서
평탄이고 $D$가 $S$ 위에서 고유이면, $D_s$가 $X_s$ 안의 카르티에
인자인 $s\in S$들의 집합(즉 $D$가 $X_s$의 점들에서 상대 카르티에
인자인 $s$들의 집합)은 $S$의 열린 부분이다. 한편 앞의 정의는 상대
카르티에 인자라는 개념이 임의의 밑변환 $S'\longrightarrow S$ 아래에서
안정하도록 정해졌다. 이제 $X/S$ 위의 상대 인자들의 집합을
$\operatorname{Div}(X/S)$라 하고, $S$ 위에서 움직이는 $S'$에 관한 반변
함자를 다음과 같이 정의하자.
\[
\operatorname{Div}_{X/S}(S')=
\operatorname{Div}(X\times_SS'/S').
\]

$X$가 $S$ 위에서 평탄이고 고유라고 가정하자. 그러면 상대 카르티에 인자의
특성화 (ii)에 따라 $\operatorname{Div}_{X/S}$를 [4], IV에서 정의한 함자
$\operatorname{Hilb}_{X/S}$의 부분함자로 볼 수 있으며, 포함 준동형
\[
\operatorname{Div}_{X/S}\longrightarrow\operatorname{Hilb}_{X/S}
\]
은 앞의 비고들에 따라 ``열린 몰입 사상들에 의해 표현 가능''하다
([5], IV, 3.13과 비교하라). [4], IV의 주요 존재 정리를 이용하면 다음을
얻는다.

\src{152}
\result{명제}{4.1} $f:X\longrightarrow S$가 사영이고 평탄하다고 하자.
그러면 함자 $\operatorname{Div}_{X/S}$는 표현 가능하며, 더 정확히는
$\operatorname{Hilb}_{X/S}$의 열린 부분준스킴으로 표현된다.

$X/S$ 위에 주어진 매우 풍부한 층 $\mathcal O_X(1)$에 대하여,
$\operatorname{Hilb}_{X/S}$를 힐베르트 다항식
$Q\in\mathbf Q[t]$에 대응하는 열린 부분준스킴
$\operatorname{Hilb}^{Q}_{X/S}$들의 합으로 분해하는 표준 분해를 사용하면,
마찬가지로 분해
\[
\operatorname{Div}_{X/S}=
\coprod_{Q\in\mathbf Q[t]}\operatorname{Div}^{Q}_{X/S}
\]
를 얻는다. 이는 $S$ 위에서 준사영인 서로소 열린 부분준스킴들의 합이다.

한편 대응 $D\longmapsto\mathcal L(D)$를 사용하면 함자적 준동형
\begin{equation}
\operatorname{Div}_{X/S}\longrightarrow\operatorname{Pic}_{X/S}
\tag{+}
\end{equation}
을 얻는다. 이제 이를 연구하고자 한다. 상당히 일반적인 조건 아래에서 이
준동형은 상대적으로 표현 가능함이 드러날 것이다([5], IV, 3). 따라서
$\operatorname{Pic}_{X/S}(S')$의 한 원소 $\xi$에서 출발하자. 표기를
간단히 하기 위하여 $S'=S$라고 가정하고, 이에 대응하는
$\operatorname{Div}_{X/S}$의 부분함자가 표현 가능함을 보이자. 먼저
$\xi$가 $X$ 위의 가역 가군 $\mathcal L$로 정의되는 경우를 생각한다.
$X$가 $S$ 위에서 고유이고 평탄하며, $S$ 위의 $X$의 기하적 섬유들이
정수적이라고 가정하자. 그러면 ([2], III, 7절)
$\mathcal O_S\xrightarrow{\sim}f_*(\mathcal O_X)$이고, 이 관계는 임의의
밑변경 $S'\longrightarrow S$ 뒤에도 성립한다. $X/S$ 위의 상대 카르티에
인자 $D$ 가운데 $\mathcal L(D)$와 $\mathcal L$이
\[
\operatorname{Pic}(X/S)=\operatorname{Pic}_{X/S}(S)
\]
의 같은 원소를 정의하는 것들, 즉 2.4에 따라 $\mathcal L(D)$와
$\mathcal L$이 $S$ 위에서 국소적으로 동형인 것들은 몫층
$f_*(\mathcal L)^*/\mathcal O_S^*$의 단면들과 일대일로 대응한다. 이
대응은 밑변경과 양립한다. 한편 인용한 곳의 “퀸네트” 유형의 일반적
고찰([6]도 보라)에 따르면, $X/S$의 고유성과 $S$ 위에서의
$\mathcal L$의 평탄성으로부터 $S$ 위의 연접 가군 $\mathcal Q$와 층의
동형사상
\[
f_*(\mathcal L)\xrightarrow{\sim}
\operatorname{Hom}_{\mathcal O_S}(\mathcal Q,\mathcal O_S)
\]
이 존재한다. 여기서 $\mathcal Q$는 유일한 동형사상을 제외하고
정의되며, 그 형성은 밑변경과 양립한다.
\footnote{인쇄본에서는 $\operatorname{Hom}$의 아래첨자에
$\mathcal O_X$를 두었다. 그러나 $\mathcal Q$는 $S$ 위의 가군이므로
그 쌍대는 $\mathcal O_S$ 위에서 취한다.} 여기서
$f_*(\mathcal L)^*$는 $f_*(\mathcal L)$의 집합들의 부분층으로서, $U$
위의 단면들은 $f^{-1}(U)$ 위의 $\mathcal L$의 단면들 가운데
$f^{-1}(U)/U$ 위의 상대 카르티에 인자를 정의하는 것들이다. 즉 이들은

\src{153}
$X_s$ 위에 영인자가 아닌 단면을 유도한다($s\in U$). 섬유 $X_s$가
정수적이라는 가정을 사용하면, 이는 유도된 단면이 각 섬유 $X_s$ 위에서
항등적으로 영이 아니라는 뜻일 뿐이다. 또는 국소 준동형
$\mathcal Q\longrightarrow\mathcal O_S$의 말로는 이 준동형들이
전사라는 뜻이다(나카야마 보조정리). 따라서
$f_*(\mathcal L)^*/\mathcal O_S^*$의 단면들의 집합은 $\mathcal Q$의
가역 몫가군들의 집합과 일대일로 대응한다. 동치로, 연접 가군
$\mathcal Q$에 대응하는 사영 다발 $\mathbf P(\mathcal Q)$의 정의에
따라([5], V, 2), 이는 $S$ 위의 $\mathbf P(\mathcal Q)$의 단면들의
집합과 일대일로 대응한다. 이 기술은 역상 형성과 양립하므로 다음 정리를
얻는다.

\result{정리}{4.3} $f:X\longrightarrow S$를 기하적 섬유가 정수적인
고유 평탄 사상이라 하고, $S$는 국소 뇌터라고 하자. 또한 $\mathcal L$을
$X$ 위의 가역 가군이라 하자. $S$ 위의 모든 $S'$에 대하여 $T(S')$를
$X\times_SS'/S'$ 위의 상대 인자 $D$ 가운데 $\mathcal L(D)$가
$S'$ 위에서 국소적으로
$\mathcal L\otimes_{\mathcal O_S}\mathcal O_{S'}$와 동형인 것들의
집합이라 하자. 즉 $\mathcal L(D)$와
$\mathcal L\otimes_{\mathcal O_S}\mathcal O_{S'}$가
$\operatorname{Pic}(X\times_SS'/S')$의 같은 원소를 정의하도록 한다.
그러면 $S$ 위의 연접 가군 $\mathcal Q$가 존재하고 유일한 동형사상을
제외하고 결정되며, 함자 $T$는 사영 다발 $\mathbf P(\mathcal Q)$로
표현된다.

\result{따름정리}{4.4} $f$가 더 나아가 사영이라고 가정하면, 함자적
준동형
$\operatorname{Div}_{X/S}\longrightarrow\operatorname{Pic}_{X/S}$은
사영 사상들로 표현된다.

실제로 $X$가 $S$ 위의 단면을 가지면(각각 국소적으로 단면을 가지면),
4.3과 2.1에 따라 앞의 준동형은 사영 다발들로(각각 국소 사영 다발들로)
표현된다. $f$가 준사영인 경우에는 $S$ 위에서 국소적으로 $X$의 유한
평탄 준단면들을 사용하는 하강 방법으로 앞의 경우에 쉽게 환원된다.

\result{비고}{4.5} 4.3의 조건 아래에서 $S$ 위의 가군 $\mathcal Q$는
일반적으로 국소 자유가 아니다. 이는 $s\in S$가 변할 때
$\mathcal Q$의 축소 섬유의 차원, 즉 $H^0(X_s,\mathcal L_s)$의 차원이
뛸 수 있다는 사실로 나타난다. 국소 뇌터 준스킴 $S$ 위의 연접 가군
$\mathcal Q$가 주어졌다고 하자. 주어진 $s\in S$에 대하여
$\mathcal Q$가 $s$에서 자유일 필요충분조건은
$\mathbf P(\mathcal Q)$가 $s$ 위의 점들에서 $S$ 위로 평탄한 것이다.
이 경우에는 그 점들에서 $S$ 위로 매끄럽기까지 하다. 여기서는
$\mathcal Q$가 위와 같이 $\mathcal L$로 정의되므로, 이는 직상
$f_*(\mathcal L)$의 형성이 “$s$의 근방에서 밑변경과 가환한다”는 뜻과도
같다. 다시 말해

\src{154}
\[
f_*(\mathcal L)_s\longrightarrow H^0(X_s,\mathcal L_s)
\]
가 전사라는 뜻이다. 예를 들어 $H^1(X_s,\mathcal L_s)=0$이면 이렇게
된다. 문제의 준스킴들이 존재한다는 전제 아래에서, 이 판정법들은 특히
보편적 상황
$\operatorname{Div}_{X/S}\longrightarrow\operatorname{Pic}_{X/S}$에
적용된다. 그리하여 이 사상이 $\operatorname{Div}_{X/S}$의 주어진 한
점에서 매끄러울 필요충분조건, 각각 충분조건을 준다.

\fsection{5}{주요 존재 정리의 증명}

3.1의 조건 아래에서 $X/S$ 위의 매우 풍부한 가군
$\mathcal O_X(1)$을 택하고, $\xi$를 이에 대응하는
$\operatorname{Pic}(X/S)$의 원소라 하자. 줄여서
\[
P(S')=\operatorname{Pic}(X\times_SS'/S')
\]
로 놓고, 먼저 표기를 간단히 하기 위하여 $X/S$가 단면을 가진다고
가정하자. $P^+(S')$를 $X\times_SS'$ 위의 가역 가군 $\mathcal L$의
류들 가운데
\[
R^if'_*(\mathcal L(n))=0\quad\text{$i>0$이고 모든 $n\geq0$에 대하여},
\qquad
f'_*(\mathcal L(n))\neq0\quad\text{모든 $n\geq0$에 대하여}
\]
를 만족하는 것들로 이루어진 $P(S')$의 부분집합이라 하자. 이 조건들은
밑변경에 대해 안정하므로 $P$의 부분함자 $P^+$를 정의하며, 이는 분명히
$\xi$에 의한 평행이동에 대해 안정하다. 세르의 “정리 A와 B”([2], III,
2)와 일반적 결과들([5], IV, 5)을 사용하면 $P$가 표현 가능할
필요충분조건은 $P^+$가 표현 가능한 것임을 쉽게 알 수 있다. 이 경우
$P^+$는 $P$를 표현하는 준스킴 $\operatorname{Pic}_{X/S}$의 열린
부분준스킴 $U$로 표현되고, 이 준스킴은 열린 부분준스킴 $U-n\xi$들의
증가 합집합이다.

줄여서 $D=\operatorname{Div}_{X/S}$로 놓고, $D^+$를 표준 사상
$D\longrightarrow P$에 의한 $P^+$의 역상이라 하자. 따라서 사상
\begin{equation}
D^+\longrightarrow P^+
\tag{+}
\end{equation}
이 있다. 또한 $D^+$는 이미 준스킴 $D=\operatorname{Div}_{X/S}$의
열린 부분준스킴 $D^+$로 표현됨을 알고 있다. 이 준스킴의 존재는 4.1이
보장한다. 4.3으로부터 사상 (+)가 매끄럽고 사영이며 전사인 사상들로
표현됨이 쉽게 따른다. 더 정확히는, 모든 점에서 영이 아닌 국소 자유
가군에 대응하는 사영 다발들로 표현된다. 실제로
$X\times_SS'$ 위의 $\mathcal L$이 이 번호의 첫머리에서와 같으면,
$f'_*(\mathcal L)$은 영이 아닌 국소 자유 가군이고 그 형성은 밑변경과
가환한다. 4.3의 표기로 $\mathcal Q$는 $f'_*(\mathcal L)$의 쌍대와
동형이다. 이제 일반 판정법([5], IV, 4.7)을 사용하여 $P^+$의 표현
가능성을 이끌어 내겠다. 인용한 곳에서 $\mathfrak C$를 충실 평탄이고

\src{155}
준콤팩트인 준스킴 사상들의 모임으로 잡는다. 이들은 [4], I B에 따라
유효 에피사상이다. 인용한 곳의 조건 (a), 즉 (+)가
$\mathfrak C$에 속하는 사상들로 표현된다는 조건은 방금 본 바와 같이
성립한다. 마찬가지로 조건 (b), 즉 함자 $P^+$가 충실 평탄이고
준콤팩트인 하강과 양립한다는 조건도 즉시 성립한다. 이제 인용한 곳의
조건 (c)를 증명하는 일만 남는다. 즉 $\mathfrak C$-표현 가능 사상
(+)로부터 준스킴 $D^+$ 위에 유도된 동치관계 $R$이
$\mathfrak C$-유효임을, 다시 말해 $R$이 유효이고
$D^+\longrightarrow D^+/R$이 $\mathfrak C$에 속함을 증명해야 한다.
이를 위해 먼저, 여러 가상 힐베르트 다항식 $Q\in\mathbf Q[t]$에 대응하는
$D^+$의 열린 부분준스킴 $D^{+Q}$가 $R$에 대해 안정함을 주목한다.
이는 $R$의 섬유들이 연결이기 때문이다. 따라서 모든 $Q$에 대하여
$D^{+Q}$ 위에 유도된 동치관계 $R^Q$가 $\mathfrak C$-유효임을 보이는
문제로 환원된다. 그런데 이제 $D^{+Q}$는 준사영이고, 한편 동치관계
$R^Q$는 사영이고 평탄하다. 그러므로 [4], III, 6.1을 적용할 수 있는
조건 아래에 있으며, 이 결과로 원하는 결론이 따른다.

$X/S$가 반드시 단면을 가지지는 않는 일반적인 경우에는 하강 기법으로
앞의 경우에 쉽게 환원한다. 또는 $P^+$의 정의에 필요한 수정을 가하여
앞의 증명을 되풀이한다.

\result{비고}{5.1} 여기서 따른 방법은 피카르 다양체를 사영적으로 구성하기
위한 마쓰사카의 방법과 본질적으로 같다. 유효 몫으로 넘어갈 수 있음을 위해
인용한 [4], III의 결과는 힐베르트 스킴의 존재 정리에서도 쉽게 도출할 수
있다(예컨대 [6]을 보라). (고전적으로 이러한 몫은 차우 좌표를 이용하여
구성했다.) $\operatorname{Pic}_{X/S}$의 열린 부분
$\operatorname{Pic}^{+}_{X/S}$의 형성과, 고려하는 가역 가군들을 정하는
약수들의 힐베르트 다항식에 따라 이를 $S$ 위에서 준사영인 열린 부분들
$\operatorname{Pic}^{+Q}_{X/S}$로 분해하는 것은 밑변경과 양립함에
주의하자. 이 사실 덕분에 하강 기법을 적용할 수 있다.

\result{비고}{5.2} $f:X\longrightarrow S$가 고유이고 평탄이며 준동형
\[
k(s)\longrightarrow H^0(X_s,\mathcal O_{X_s})\qquad(s\in S)
\]
들이 동형사상일 때마다 $\operatorname{Pic}_{X/S}$가 존재할 가능성을
배제할 수 없다. (이 마지막 조건은 이때
$\mathcal O_S\xrightarrow{\sim}f_*(\mathcal O_X)$이고 이 관계가 임의의
밑변경 $S'\longrightarrow S$ 뒤에도 유지된다는 뜻이기도 하다.) 적어도
$f$가 더 나아가 사영인 해석공간의 경우에는, [5], IX, 3.1에서 설명한
다른 방법(내가 틀리지 않았다면 차우의 방법)으로 이 사실을
증명한다.\footnote{1962년 정오표는 이 비고의 첫머리에서 제시한 존재
추측이 거짓이지만, 멈퍼드의 방법으로 그보다 약간 약한 정리를 얻을 수
있다고 지적한다.} 그 방법에서는 약수 스킴의 한 열린 부분 위에서 고유이고
평탄인 동치관계로 몫을 취하는 대신, $X$를 $\mathbf P^r_S$ 안에
몰입시키는 몰입 사상들의 스킴에서 사영군으로 몫을 취한다.

\src{156}
이 방법은 사영군으로 몫을 취하는 것에 관한 멈퍼드의 결과 [8]을 이용하면
스킴의 경우에도 아마 적용할 수 있을 것이다. 현재로서는 $X$가 $S$ 위에서
``많은 국소 단면''을 갖는 경우에만 증명이 쓰여 있다. 예컨대 $X$가
잉여체가 대수적으로 닫힌 완비 국소환 위에서 분리 가능한 경우가
그렇다.\footnote{인쇄본에는 ``잉여체가 대수적으로 닫힌''이라는 말이
빠져 있다. 정오표는 멈퍼드의 반례 때문에 이 제한이 반드시 필요하다고
명시한다.} 원칙적으로 이 방법은 더 일반적으로 적용될 것이다. 실제로
피카르 준스킴이 분리되지 않는 경우에도 그 존재를 보이므로, 그런 경우에는
첫 번째 방법이 필연적으로 실패한다. (기술적인 어려움은 $f$의 기하적
섬유들이 더 이상 정역이 아닐 때 4.3에서 고려한 함자가 사영 다발
$\mathbf P(\mathcal Q)$ 자체로 표현되는 것이 아니라 그 안의 열린 부분으로
표현된다는 데서 생긴다. 이 때문에 평탄이지만 고유이지 않은 동치관계로
몫을 취하는 까다로운 문제가 나타난다.)

\result{비고}{5.3} 여기서 제시한 증명은 곡선 또는 곡선족의 야코비안을
미리 구성하는 일에도, 아벨 다양체나 아벨 스킴의 이론에도 의존하지 않는다.
따라서 A. 베유가 개략적으로 제시한 길을 따르는 랑의 책 [7]이나 슈발레의
강연 [1] 같은 전통적인 설명과 본질적으로 다르다. 대수적으로 닫힌 체
(가령 복소수체) 위의 비특이 곡선의 야코비안인 경우에도, 여기서 제시한
야코비안의 구성은 1번에서 정의로 삼은 매우 강한 성질들을 보장하는 것으로
알려진 유일한 구성이다. (이는 본질적으로 슈발레의 성질들이지만, 멱영원을
갖는 ``매개변수 다양체''까지 고려한 것이다.) 피카르 스킴의 구성은 아벨
다양체 이론 뒤에 와야 하는 것이 아니라 그에 앞서야 한다. 이는 일반적으로
피카르 스킴이 아벨 스킴도 아니고 아벨 스킴으로 환원되지도 않는다는
사실에서 선험적으로 분명하다. 대수적으로 닫힌 체 위의 특이 곡선에서 이미
이를 볼 수 있다. 그 경우에는 아벨 다양체가 아닌 로젠리히트의 ``일반화
야코비안''이 나타난다. 더 나아가 피카르 스킴의 일반 이론을 갖추고 나면
아벨 다양체 이론, 더 일반적으로는 아벨 스킴 이론이 훨씬 간단해진다.
특히 아벨 스킴의 쌍대성 이론, 그중에서도 카르티에형 결과들은 이로부터
거의 형식적인 귀결이 된다(예컨대 [10]을 보라).

\src{157}
\result{비고}{5.4} 특이 곡선으로 퇴화하는 곡선의 야코비안에 관한
이구사의 ``양립성 원리''는, 반드시 $S$ 위에서 매끄럽지는 않은 곡선을
섬유로 갖는 상대 스킴 $X/S$의 피카르 스킴에 대한 존재 정리로 이해해야만
제대로 이해할 수 있다. 따라서 특수화된 곡선이 정역일 때, 즉 고전적
용어로 기약이고 중복도가 1일 때 이는 주요 존재 정리 3.1의 특수한 경우다.
현재로서는 특수 곡선이 가약인 경우는 알려진 존재 정리들로 다룰 수 없다.
각 성분의 중복도가 1인 경우, 즉 특수 곡선이 잉여체 위에서 분리 가능한
경우에도 마찬가지다. 다만 잉여체가 대수적으로 닫힌 완비 이산 값매김환
위에서는 예외이며, 5.2와 비교하라. 종수 $g$인 곡선들의 모듈라이 스킴에
대한 베일리--사타케 ``콤팩트화''를 스킴 이론 안에서 기하적으로 구성할
때 이 존재 문제는 분명 제기될 것이다. (이 콤팩트화는 현재 이구사의
연구 덕분에 $g=1$인 경우에만 알려져 있다.)

\fsection{6}{상대 존재 정리}

우리는 여기서 어떤 피카르 스킴들의 존재가 다른 피카르 스킴들의 존재를
함의하는 몇 가지 유용한 경우를 개략적으로 살펴본다. 이로써 주요 정리
3.1에서 여러 다른 존재 정리를 도출할 수 있다.

\result{명제}{6.1} $f:X\longrightarrow S$를 평탄 사영 사상이라 하자. 스타인
인수분해 $f=f''f'$에서 사상 $f':X\longrightarrow S'$은 평탄이고 그 기하적
섬유들이 정역이며(따라서 3.1의 가정들을 만족한다), 유한 사상
$f'':S'\longrightarrow S$은 평탄하다고 하자. 그러면
$\operatorname{Pic}_{X/S}$가 존재하고, [4], II, C 2에서 도입한 표기를 쓰면
표준 동형사상
\[
\operatorname{Pic}_{X/S}\xrightarrow{\sim}
\Pi_{S'/S}\operatorname{Pic}_{X/S'}
\]
이 있다.

이를 보이기 위해 먼저 함자들의 동형사상
\[
\operatorname{Pic}_{X/S}\xrightarrow{\sim}
\Pi_{S'/S}\operatorname{Pic}_{X/S'}
\]
을 확립하고, $\operatorname{Pic}_{X/S'}$가 표현 가능함을 함의하는 3.1을
사용한다. 이어서 제3절에서 지적한 $\operatorname{Pic}_{X/S'}$의 구조, 즉
$S$ 위의 $\operatorname{Pic}_{X/S'}$의 한 섬유의 모든 유한 부분집합이 어떤
아핀 열린집합에 포함된다는 사실을 사용하여
$\Pi_{S'/S}\operatorname{Pic}_{X/S'}$의 존재를 얻는다.

\src{158}
예를 들어 $X$가 3.1의 조건들을 만족하는 $S$ 위의 스킴 $X_i$들의 서로소
합인 스킴이면, 명제 6.1은 자명한 명제
\[
\operatorname{Pic}_{X/S}\xrightarrow{\sim}
\prod_i\operatorname{Pic}_{X_i/S}
\]
로 환원된다.

\result{따름정리}{6.2} $f:X\longrightarrow S$를 기하적 섬유들이 국소 정역인
사영 평탄 사상이라 하자(예를 들어 사영이고 정규인 사상). 그러면
$\operatorname{Pic}_{X/S}$가 존재한다.

더욱이 이 경우 $S'$은 $S$의 에탈 피복이다. 이는 $f$가 분리 가능할 때, 즉
평탄이고 그 기하적 섬유들이 축소일 때마다 참이다. 또한
$\operatorname{Pic}_{X/S'}$의 유사한 구조 덕분에, 제3절에서
$\operatorname{Pic}_{X/S}$에 대해 진술한 구조정리가 여전히 성립함을 확인할
수 있다.

한편 하강 절차는 상대 존재정리를 준다. 그 적용 범위는 [4], I, A (c)에서
제기한 비평탄 하강 문제들의 해결에 달려 있는데, 여기서는 다음 특수한
경우만 명시하겠다.

\result{명제}{6.3} $f:X\longrightarrow S$를 고유 사상이라 하고, $X_1$과
$X_2$를 $S$ 위에서 평탄인 $X$의 두 부분준스킴이라 하자. 이들이 두 연접
아이디얼 $\mathcal I_1$과 $\mathcal I_2$로 정의되고
$\mathcal I_1\cap\mathcal I_2=(0)$이며
$\mathcal O_X/(\mathcal I_1+\mathcal I_2)$가 $S$ 위에서 평탄이라고 하자.
즉 $X$의 부분준스킴 가운데 $X_1$과 $X_2$의 상한은 $X$이고, 그 하한 $Z$는
$S$ 위에서 평탄이라고 하자.\footnote{마지막 글자는 인쇄본에서 ``$X$''로
되어 있으나, 바로 앞 조건은 $\mathcal O_X/(\mathcal I_1+\mathcal I_2)$가
$S$ 위에서 평탄하다는 것이므로 $S$ 위에서 평탄인 것은 $Z$이다.} 또한 모든
$s\in S$에 대하여 준동형사상
\[
k(s)\longrightarrow H^0(X_{i,s},\mathcal O_{X_{i,s}}),\qquad i=1,2
\]
가 전단사라고 가정하자. 그러면 함자들의 자연 준동형사상
\[
\operatorname{Pic}_{X/S}\longrightarrow
\operatorname{Pic}_{X_1/S}\times\operatorname{Pic}_{X_2/S}
\]
은 아핀 사상으로 표현 가능하다. 따라서 $\operatorname{Pic}_{X_1/S}$와
$\operatorname{Pic}_{X_2/S}$가 존재하면 $\operatorname{Pic}_{X/S}$도
존재하고, 표준 사상
\[
\operatorname{Pic}_{X/S}\longrightarrow
\operatorname{Pic}_{X_1/S}\times\operatorname{Pic}_{X_2/S}
\]
은 아핀이다.

충실 평탄 하강에 의해 $Z$가 $S$ 위의 단면을 갖는 경우로 환원할 수 있다.
이 단면은 $X$, $X_1$, $X_2$의 $S$ 위 단면들을 정하므로 2.5에서 설명한
대로 지금 고려하는 구조들에서 자기동형사상들을 제거할 수 있다. 이제 증명은
$X$ 위의 ``강체화된'' 가역 가군 $\mathcal L$을 주는 것이 삼중항

\src{159}
$(\mathcal L_1,\mathcal L_2,u)$를 주는 것과 동치임을 지적하는 데 있다.
여기서 $\mathcal L_i$는 $X_i$ 위의 ``강체화된'' 가역 가군이고, $u$는
강체화들과 양립하는 $\mathcal L_1|Z$에서 $\mathcal L_2|Z$로의
동형사상이다. $\mathcal L_1$, $\mathcal L_2$가 고정되었을 때 $u$를 주는
것이 $S$ 위에서 아핀인 적당한 $S$ 위 스킴의 한 단면으로 표현됨을 확인하기만
하면 되는데, 이는 쉽다. 6.3에서 곧바로 다음을 얻는다.

\result{따름정리}{6.4} $X$를 체 $k$ 위의 고유이고 분리 가능한 스킴이라
하고, $X_i$들을 $X$의 기약 성분들이라 하자. 모든
$\operatorname{Pic}_{X_i/k}$가 존재하면 $\operatorname{Pic}_{X/k}$도
존재하고, 표준 사상
\[
\operatorname{Pic}_{X/k}\longrightarrow\prod_i\operatorname{Pic}_{X_i/k}
\]
은 아핀이다.

이를 6.2와 결합하면, 예를 들어 $X$가 체 $k$ 위의 사영이고 분리 가능한
스킴일 때마다 $\operatorname{Pic}_{X/k}$가 존재함을 알 수 있다. $X$가 더는
$k$ 위에서 분리 가능하지 않을 때에도 Oort [11]의 논증을 사용하면 환원
결과를 얻는다. 이 방법은 임의의 밑스킴 위에서도 적용된다(예를 들어 다음
강연에서 3.3에 예고한 유한성 결과를 증명할 때 유용하다). 지나치게 기술적인
진술을 피하기 위해 기초체 위의 경우로 한정하겠다.

\result{명제}{6.5} $X$를 체 $k$ 위의 고유 스킴이라 하고, $X_0$를 $X$와
같은 바탕집합을 갖는 부분스킴이라 하자(따라서 $X$ 위의 한 멱영 아이디얼로
정의된다). 그러면 함자적 사상
\[
\operatorname{Pic}_{X/k}\longrightarrow\operatorname{Pic}_{X_0/k}
\]
은 아핀 사상으로 표현 가능하다. 특히 $\operatorname{Pic}_{X_0/k}$가
존재하면 $\operatorname{Pic}_{X/k}$도 존재하고, 사상
$\operatorname{Pic}_{X/k}\longrightarrow\operatorname{Pic}_{X_0/k}$는
아핀이다.

이를 6.4와 결합하면 곧바로 다음을 얻는다.

\result{따름정리}{6.6} $X$를 체 $k$ 위의 사영 준스킴이라 하자. 그러면
$\operatorname{Pic}_{X/k}$가 존재한다.

\result{비고}{6.6} 체 $k$ 위의 모든 고유 스킴 $X$에 대하여
$\operatorname{Pic}_{X/k}$가 존재한다는 것은 지극히 그럴듯하다.\footnote{
1962년 정오표는 Murre가 이 문제를 막 긍정적으로 해결했다고 알린다.} 앞의
결과들은 이 문제를 $k$가 대수적으로 닫혀 있고 $X$가 정역인 경우로
환원한다. 그러면 차우 보조정리에 의해 $k$ 위에서 사영인 정역 스킴 $X'$과
우세 사상 $g:X'\longrightarrow X$가 존재한다. 따라서 대응하는 함자적 사상

\src{160}
\[
\operatorname{Pic}_{X/k}\longrightarrow\operatorname{Pic}_{X'/k}
\]
이 표현 가능함을(그리고 아마도 아핀 사상으로 표현 가능함을) 보이면 충분할
것이다. 실제로 $\operatorname{Pic}_{X'/k}$가 표현 가능함은 이미 알고 있기
때문이다. 이는 현재 해결되지 않은 비평탄 하강 문제들을 제기한다. $X$가
대수적으로 닫힌 체 $k$ 위의 고유 정역 스킴일 때, 함자
$\operatorname{Pic}_{X/k}$를 축소 준스킴들에 제한하면 표현 가능한 함자를
얻는다는 점에 유의하자. $X$가 정규인 경우에는 Chevalley [1]가 이를
보였으며, 일반적인 경우에는 Seshadri [12]가 하강 방법으로 보였다. 그러나
우리 표기에서 이 저자들이 구성한 스킴은 $\operatorname{Pic}_{X/k}$가 아니라
그에 대응하는 축소 스킴
$(\operatorname{Pic}_{X/k})_{\mathrm{réd}}$이다.

\result{비고}{6.7} 더 일반적으로 $f:X'\longrightarrow X$를 $k$ 위의 고유
준스킴들의 전사 사상이라 하자. 그러면 비평탄 하강을 고찰하면
\[
\operatorname{Pic}_{X/k}\longrightarrow\operatorname{Pic}_{X'/k}
\]
가 아핀 사상일 것이라고 추측하게 된다.\footnote{1962년 정오표는 이 문제가
긍정적으로 해결되었다고 알리며 다음 강연의 주석 마지막 문단을 참조하도록
한다.} 이는 특히 항등원의 연결 성분들로 나눈 뒤 대응하는 네롱--세베리
군들의 준동형사상이 꼬임을 법으로 단사임을 함의할 것이다. $X$와 $X'$이
비특이이면 교차이론으로 이를 확인할 수 있다. 그 밖의 모든 경우에는 답이
알려져 있지 않은 듯하다.

\result{비고}{6.8} 생각과 달리 멱영원을 갖는 대수적 스킴들의 피카르
스킴을 고찰하는 일은 여러 문제에서 유용하고, 심지어 필수적이다. 예를 들어
$X$가 $k$ 위의 사영 스킴, 이를테면 매끄러운 스킴이고,\footnote{인쇄본에는
``정칙''이라고 되어 있으나 정오표는 ``$k$ 위에서 매끄러운''으로 바로잡도록
한다.} $Y$가 한 초평면 단면이면, $Y$의 모든 차수의 ``무한소 근방'' $X_n$과
피카르 스킴 $\operatorname{Pic}_{X_n/k}$를 고려해야 한다. $X$가 차원
$\geq4$(각각 $\geq3$)인 기약 스킴이면 표준 사상
\[
\operatorname{Pic}_{X/k}\longrightarrow\operatorname{Pic}_{X_n/k}
\qquad(n\text{이 충분히 클 때})
\]
은 동형사상이다(각각 네롱--세베리 꼬임 부분군들의 역상들 사이에
동형사상을 유도한다). 이 결과는 다음 강연에서 피카르 스킴을 정성적으로
연구할 때 유용할 것이다. 마찬가지로 멱영원을 갖는 몇몇 곡선들의 피카르
스킴을 고찰하고 형식기하학의 기본정리들 [3]을 사용하면, 동표수인 경우에
Mumford의 한 추측을 확인할 수 있다. 즉 차원 2인 모든 완비 정규 뇌터 국소환
$A$에 대하여, $A$의 약수류군을 잉여체 $k$ 위의 한 대수군 $G$의 $k$-유리점
집합으로 볼 수 있다.

\src{161}
더욱이 $A$ 안에서 계수체 하나를 택하면 $G$는 표준적으로 결정된다. $A$의
차원이 임의일 때에는 위의 $G$와 같은 역할을 하는 $k$ 위의 프로대수군이
존재할 법하다. $\operatorname{Spec}(A)$의 ``특이점을 해소''할 수 있는
경우, 이 군은 $k$ 위의 적당한 사영 스킴들(멱영원을 갖는)의 피카르 스킴들의
사영극한으로 구성된다.

\section*{참고문헌}
\addcontentsline{toc}{section}{참고문헌}

\begin{description}
\item[{[1]}] CHEVALLEY (Claude). -- Sur la théorie de Picard,
\emph{Amer. J. of Math.}, 제82권, 1960, 435--490쪽.

\item[{[2]}] GROTHENDIECK (A.) et DIEUDONNÉ (J.). -- Éléments de géométrie
algébrique, 제I장 이하. -- Paris, Presses universitaires de France,
1960--1961 (Institut des Hautes Études Scientifiques, Publications
mathématiques 4, 8, \ldots).

\item[{[3]}] GROTHENDIECK (Alexander). -- Géométrie formelle et géométrie
algébrique, \emph{Séminaire Bourbaki}, 제11권, 1958/59,
제182호, 28쪽.

\item[{[4]}] GROTHENDIECK (Alexander). -- Techniques de descente et théorèmes
d'existence en géométrie algébrique, I--IV, \emph{Séminaire Bourbaki}, 제12권,
1959/60, 제190호, 29쪽 및 제195호, 22쪽; 제13권, 1960/61, 제212호, 20쪽 및
제221호, 28쪽.

\item[{[5]}] GROTHENDIECK (Alexander). -- Techniques de construction en
géométrie analytique, I--X, \emph{Séminaire Cartan}, 제13권, 1960/61:
Déformation des structures analytiques complexes, 제7--17호.

\item[{[6]}] GROTHENDIECK (Alexander). -- Séminaire de géométrie algébrique,
Lichtenbaum 집필, Harvard University, 1961 (출간 예정).

\item[{[7]}] LANG (Serge). -- \emph{Abelian varieties}. -- New York,
Interscience Publishers, 1959 (Interscience Tracts in pure and applied
Mathematics, 7).

\item[{[8]}] MUMFORD (D.). -- An elementary theorem in geometric invariant
theory, \emph{Bull. Amer. math. Soc.}, 제67권, 1961, 483--486쪽.

\item[{[9]}] MUMFORD (D.). -- The topology of normal singularities of an
algebraic surface and a criterion for simplicity. -- Paris, Presses
universitaires de France, 1961 (Institut des Hautes Études Scientifiques,
Publications mathématiques, 9).

\item[{[10]}] MUMFORD (D.) and TATE (J.). -- Séminaire de géométrie
algébrique, Harvard University, 1962년 봄학기 (출간 예정).

\item[{[11]}] OORT (F.). -- Sur le schéma de Picard,
\emph{Bull. Soc. math. France}, 제90권, 1962 (출간 예정).

\item[{[12]}] SESHADRI (C. S.). -- \emph{Annali di Matematica}에 게재 예정인
박사학위논문.
\end{description}

% FGA Canon print-preservation apparatus: historical integration R25.
\par\smallskip
\begingroup\footnotesize
\noindent\textit{FGA-CANON-ERR-0043. 인쇄본의 3.3 참조는 C109에 계속 기록되어 있으며, FGA-CANON-DISP-0043이 3.2로 바로잡은 참조를 규율한다.}\par
\endgroup
